\documentclass[11pt,a4paper]{article}
\title{Delhi-NCR Hourly PM2.5 Spatiotemporal Forecasting}
\author{Research Team}
\date{October 2026}

\begin{document}
\maketitle

\begin{abstract}
This study forecasts hourly PM2.5 concentrations across approximately 70 monitoring stations in the Delhi-NCR region. The core contribution is comparing static distance-based graphs with wind-aware dynamic graphs within a Graph WaveNet framework. All experiments use chronological splits (2019--2022 train, 2023 validation, 2024 test), masked evaluation on observed values only, and normalization statistics fitted exclusively on the training period.
\end{abstract}

\section{Introduction}

Air quality forecasting in Delhi-NCR remains a critical public health challenge. We address spatiotemporal forecasting of hourly PM2.5 across approximately 70 CPCB/CAAQMS stations.

\section{Methodology}

\subsection{Data and Frozen Dataset}
Data is frozen as a single \texttt{.npz} file before modeling begins: \texttt{delhi\_ncr\_pm25\_2019\_2024.npz}. Arrays include input \texttt{X}, target \texttt{y}, masks \texttt{mask\_X} and \texttt{mask\_y}, station IDs, coordinates, and feature names. Normalization uses \texttt{log1p(PM2.5)} followed by z-score, fitted on train data only. Any correlation-based adjacency is also computed on train data only.

\subsection{Graph Construction}
Four graph types are implemented:
\begin{itemize}
    \item Distance: Gaussian kernel on station coordinates.
    \item Correlation: Pairwise PM2.5 correlation (train only, thresholded).
    \item Wind-aware directed: Edge weights based on wind alignment.
    \item Learned adaptive: Self-adaptive adjacency \texttt{A = softmax(ReLU(E_1 E_2^T))}.
\end{itemize}

\subsection{Models}
Baselines follow the exact order in the master spec: persistence, historical average, LightGBM, univariate LSTM, multivariate LSTM, plain transformer, STGCN, and Graph WaveNet. A minimal GCN+GRU (~150 lines) is also implemented for message-passing understanding.

\subsection{Training Protocol}
Masked MAE loss. Adam optimizer. Learning rate $2 \times 10^{-4}$ to $3 \times 10^{-3}$. Batch size 64 (reduced if OOM). Dropout 0.3. Early stopping (patience 15) on validation MAE. Gradient clipping at 1.0. Three random seeds; results reported as mean $\pm$ std.

\section{Results}

Results are logged continuously in \texttt{EXPERIMENTS.md}. Final evaluation tables will be completed after full real-data runs.

\section{Conclusion}

The pipeline is fully reproducible. Baselines are logged. The dataset freeze mechanism prevents leakage. The main claim (wind-aware vs static distance graph) is structured for honest reporting regardless of outcome.

\end{document}
